% Generated draft: requires mathematical and source review.
\documentclass[11pt]{article}
\usepackage[margin=25mm]{geometry}
\usepackage{fontspec}
\setmainfont{Latin Modern Roman}
\usepackage{amsmath,amssymb,mathtools,unicode-math}
\setmathfont{Latin Modern Math}
\usepackage{enumitem,needspace,xurl,hyperref}
\setlength{\parindent}{0pt}
\setlength{\parskip}{5pt}
\setlength{\emergencystretch}{3em}
\newcommand{\N}{\mathbb{N}}
\newcommand{\Z}{\mathbb{Z}}
\newcommand{\Q}{\mathbb{Q}}
\newcommand{\R}{\mathbb{R}}
\newcommand{\C}{\mathbb{C}}
\newcommand{\F}{\mathbb{F}}
\newcommand{\A}{\mathbb{A}}
\newcommand{\PP}{\mathbb{P}}
\newcommand{\E}{\mathbb{E}}
\newcommand{\eps}{\varepsilon}
\newcommand{\dd}{\,\mathrm{d}}
\begin{document}
\section*{UnsolvedMath ID: 800882}
\textbf{Source problem code:} catalogue entry
\subsection{Definitions and mathematical statement}
Let \(n\geq 1\), and let \(\Omega\subset\C^n\) denote the intended domain. A function is plurisubharmonic (psh) if it is upper semicontinuous and its restriction to each complex line is subharmonic wherever defined. For a continuous psh function \(v\), write \((dd^c v)^n\) for its complex Monge--Amp\`ere measure in the locally bounded sense, using one consistent normalization of \(d^c\). For nonnegative Borel measures, \(\mu\leq\nu\) means \(\mu(E)\leq\nu(E)\) for every Borel set \(E\subset\Omega\).

The question attributed to Ko\l{}odziej asks whether the existence of a continuous psh subsolution with zero boundary values guarantees a continuous psh solution with the same boundary values. Interpreting the stated continuity and boundary values as continuous extension to \(\overline\Omega\), the proposed implication is: for every nonnegative Borel measure \(\mu\) on \(\Omega\),
\[
\begin{gathered}
\exists u\in\operatorname{PSH}(\Omega)\cap C(\overline\Omega):\qquad
u:= (dd^c u)^n\geq\mu,\qquad u|_{\partial\Omega}=0
\end{gathered}
\quad\Longrightarrow\quad
\begin{gathered}
\exists\varphi\in\operatorname{PSH}(\Omega)\cap C(\overline\Omega):\qquad
(dd^c\varphi)^n=\mu,\qquad\varphi|_{\partial\Omega}=0.
\end{gathered}
\]
The measure inequality and equality are required throughout \(\Omega\).

The supplied record does not specify the intended class of domains, including boundedness, boundary regularity, or pseudoconvexity assumptions. It also does not explicitly specify the measure class, normalization of \(d^c\), or boundary-trace convention. The displayed formulation makes the usual nonnegative-measure and continuous-extension interpretations explicit; it does not resolve these missing source qualifications or assert the implication for every domain in \(\C^n\).
\subsection{Short English statement}
The question attributed to Ko\l{}odziej asks whether a continuous plurisubharmonic subsolution whose Monge--Amp\`ere measure dominates \(\mu\), with zero boundary values, ensures a continuous plurisubharmonic solution with Monge--Amp\`ere measure exactly \(\mu\) and the same boundary values. The intended domain hypotheses and precise boundary-continuity convention are missing from the supplied record.
\subsection{Sources}
\begin{itemize}
\item Local UnsolvedMath registry, permanent ID \(800882\), source code \texttt{catalogue entry}, title: \emph{Open Problems in Pluripotential Theory --- Question 13}. The supplied registry labels the problem open; this status has not been independently verified.
\item Dinew, Guedj, and Zeriahi, \emph{Open Problems in Pluripotential Theory} (2015), Question 13, attributed in the supplied statement to Ko\l{}odziej. Supplied source URL: \url{https://arxiv.org/abs/1511.00705}.
\item Supplied provenance records access on 2026-10-16 and extraction by \texttt{arxiv\_latex}. These are supplied provenance details, not a claim of independent source consultation or verification.
\end{itemize}
\end{document}
